\documentclass[12pt]{article}
\usepackage[margin=1in]{geometry}
\usepackage{booktabs}
\usepackage{amsmath}
\usepackage{natbib}
\usepackage[hidelinks]{hyperref}
\usepackage{setspace}
\onehalfspacing

\title{Benchmark choice and the apparent long-horizon\\
predictive content of financial conditions indices}

\author{Franz X. Mohr\thanks{Replication code and data are available in the
\texttt{fincond} R package, \url{https://github.com/franzmohr/fincond}.}}

\date{\today}

\begin{document}

\maketitle

\begin{abstract}
\noindent Financial conditions indices are usually evaluated against an
autoregressive benchmark. For a near-unit-root target in a short sample that
benchmark is weak, and it becomes weaker as the horizon lengthens. Estimating
the time-varying-parameter FAVAR of \citet{koop2014} recursively for the euro
area and ten member countries, we find that a random walk beats the own-lag
regression in ten of eleven regions at every horizon, by a median of 33\% of
root mean squared error at two years. Against the better of a random walk and a
recursive mean, the index improves unemployment forecasts one and two quarters
ahead in ten and nine of the eleven regions, but at two years it loses in all
eleven. Against the own-lag regression alone the same forecasts appear to win at
every horizon. Long-horizon predictive content reported against autoregressive
benchmarks should be treated as a statement about the benchmark.

\vspace{1em}
\noindent \textbf{Keywords:} financial conditions index; forecast evaluation;
benchmark; dynamic factor model; euro area

\noindent \textbf{JEL:} C53, E44, E37
\end{abstract}

\section{Introduction}

Financial conditions indices compress many financial series into one number,
and the case for them rests on whether that number helps forecast real
activity. The standard evaluation regresses a target variable at $t+h$ on its
own lag and the index, and compares the resulting forecast errors with those of
the same regression without it \citep{koop2014, hatzius2010}.

The own-lag regression is a natural benchmark and a weak one. An unemployment
rate is close to a unit root. Estimating its autoregressive coefficient on
eighty quarters shrinks the estimate below one and pulls the forecast toward a
sample mean that the level has drifted away from, and the induced bias
compounds with the horizon. A random walk, which imposes the unit root rather
than estimating it, does not have this problem.

This note asks what the benchmark is worth. Using one index, one model and one
dataset for eleven euro-area regions, we compare the index against three simple
benchmarks: the recursive mean of the target, a random walk, and the own-lag
regression. The random walk is the strongest of the three in ten of the eleven
regions at every horizon, and its advantage over the own-lag regression widens
from a median of 4\% of root mean squared error at one quarter to 33\% at two
years.

The consequences for what the index is judged to be worth are large. Against
the better of the two competent benchmarks, the index improves on unemployment
forecasts at one and two quarters --- in ten and nine of eleven regions, with
Clark--West statistics clearing the one-sided five per cent value in eleven and
ten --- and loses in all eleven regions at two years. Against the own-lag
regression alone, the same forecasts appear to win at one, two, four and eight
quarters in eleven, eleven, eleven and ten regions respectively.

We read this as a positive result about short horizons and a negative one about
long horizons. Financial conditions carry genuine information about
unemployment one and two quarters out, across a panel of countries and against
a benchmark that is hard to beat. Reported evidence of predictive content at
horizons of a year or more, when it rests on an autoregressive benchmark, may be
measuring the benchmark's deterioration rather than the index's information.

\section{Method}

Let $y_t$ collect output growth, inflation and the unemployment rate, and $x_t$
the financial series. The index is the first factor $f_t$ of the FAVAR of
\citet{koop2014},
\begin{align}
x_t &= \Lambda^y_t y_t + \Lambda^f_t f_t + e_t, \qquad e_t \sim N(0, V_t), \\
z_t &= \sum\nolimits_{i=1}^{p} B_{it} z_{t-i} + u_t, \qquad u_t \sim N(0, Q_t),
\end{align}
with $z_t = (y_t', f_t')'$. Loadings and transition coefficients follow random
walks; the error covariances and the drift are governed by forgetting factors
rather than estimated variances, so the model follows from two Kalman passes
rather than from simulation. We use the authors' settings throughout ($p = 4$,
one factor, forgetting factors $0.96$, $0.96$, $0.99$, $0.99$) and fix the sign
of the factor on the economic sentiment indicator.

The evaluation is recursive. At each origin $t$ the model is re-estimated on
data through $t$ only, giving the index path a user would have had at that
date. Four forecasts of the target $w_{t+h}$ are then formed, each using data
through $t$ only:
\begin{align}
\text{mean:} \quad & \widehat{w} = \textstyle\frac{1}{t-h}\sum_{s} w_{s}, &
\text{random walk:} \quad & \widehat{w} = u_t, \\
\text{own lag:} \quad & \widehat{w} = \hat\alpha + \hat\beta u_t, &
\text{index:} \quad & \widehat{w} = \hat\alpha + \hat\beta u_t + \hat\gamma f_t .
\label{eq:pred}
\end{align}
Nothing estimated at $t$ uses data after $t$, the regression coefficients
included. All three benchmarks are nested in the index regression: the mean
sets $\beta = \gamma = 0$, the random walk fixes $\alpha = 0$, $\beta = 1$,
$\gamma = 0$.

We take the target to be the average of the variable over $t+1, \dots, t+h$
rather than its value in the single quarter $t+h$. The two coincide at $h = 1$.
Section~\ref{sec:robust} reports the single-quarter target, which is less
favourable to the index at every horizon; using the path target makes the
negative result below conservative.

We report the ratio of the index regression's root mean squared forecast error
to that of the best-performing benchmark, and the statistic of
\citet{clark2007}. Because the models are nested, an unadjusted
Diebold--Mariano test is undersized: under the null the larger model estimates
a coefficient that is truly zero, and the noise in doing so inflates its
squared error. The Clark--West adjustment removes that penalty. Overlapping
$h$-step forecasts are handled with a Bartlett kernel of $h-1$ lags, and we use
the one-sided normal critical value of $1.645$.

\section{Data}

We use EA-MD-QD \citep{barigozzi2024}, which provides quarterly series for the
euro area aggregate and ten member countries on a common definition from
2000\,Q1. The financial block has twenty-two series per region: equity and
house prices, effective and dollar exchange rates, short and long interest
rates and three spreads, sectoral balance sheets for households, firms and
banks, four confidence balances and two global risk measures. Austria and
Germany carry four national series in addition. Forecast origins run from
2006\,Q1, giving 72 to 79 forecasts per region and horizon.

\section{The benchmark}

Table~\ref{tab:bench} reports the ratio of the own-lag regression's root mean
squared error to the random walk's. The regression is worse everywhere except
Austria, and the gap widens with the horizon: a median of 1.04 at one quarter
and 1.33 at two years, reaching 2.62 for Ireland. The mechanism is the one
described in the introduction, and Austria is the exception that confirms it,
being the one region whose unemployment rate mean-reverts over this sample
rather than drifting.

\begin{table}[htbp]
\centering
\caption{Own-lag regression relative to a random walk, RMSFE ratio}
\label{tab:bench}
\begin{tabular}{lrrrr}
\toprule
& $h{=}1$ & $h{=}2$ & $h{=}4$ & $h{=}8$ \\
\midrule
Euro area   & 1.03 & 1.09 & 1.22 & 1.34 \\
Austria     & 0.97 & 0.95 & 0.89 & 0.83 \\
Belgium     & 1.02 & 1.04 & 1.11 & 1.25 \\
Germany     & 1.05 & 1.11 & 1.26 & 1.66 \\
Greece      & 1.05 & 1.18 & 1.52 & 2.07 \\
Spain       & 1.05 & 1.11 & 1.22 & 1.61 \\
France      & 1.03 & 1.07 & 1.14 & 1.25 \\
Ireland     & 1.12 & 1.42 & 2.00 & 2.62 \\
Italy       & 1.04 & 1.09 & 1.19 & 1.33 \\
Netherlands & 1.04 & 1.08 & 1.14 & 1.28 \\
Portugal    & 1.04 & 1.10 & 1.20 & 1.32 \\
\midrule
Median      & 1.04 & 1.09 & 1.20 & 1.33 \\
\bottomrule
\end{tabular}

\vspace{0.5em}
\begin{minipage}{0.9\textwidth}
\footnotesize \textit{Notes.} Values above one mean the random walk forecasts
the path of the unemployment rate more accurately than a recursively estimated
own-lag regression. Origins from 2006\,Q1.
\end{minipage}
\end{table}

\section{Results}

Table~\ref{tab:main} scores the index against whichever of the three
benchmarks did best in each cell. The index wins at one quarter in ten of
eleven regions and at two quarters in nine, with Clark--West statistics
clearing the critical value in eleven and ten. At four quarters it is close to
a coin flip. At two years it loses in every region, by 32\% of root mean
squared error on average, and the Clark--West statistic clears in only three.

\begin{table}[htbp]
\centering
\caption{Index against the best of three benchmarks, unemployment}
\label{tab:main}
\begin{tabular}{lrrrrcrrrr}
\toprule
& \multicolumn{4}{c}{RMSFE ratio} & & \multicolumn{4}{c}{Clark--West} \\
\cmidrule(lr){2-5} \cmidrule(lr){7-10}
& $h{=}1$ & $h{=}2$ & $h{=}4$ & $h{=}8$ & & $h{=}1$ & $h{=}2$ & $h{=}4$ & $h{=}8$ \\
\midrule
Euro area   & 0.770 & 0.831 & 0.979 & 1.206 & & 3.31 & 2.76 & 1.77 & 0.64 \\
Austria     & 0.958 & 0.931 & 0.915 & 1.086 & & 3.03 & 3.14 & 3.38 & 1.72 \\
Belgium     & 0.957 & 0.915 & 0.952 & 1.204 & & 2.32 & 2.96 & 2.44 & 1.88 \\
Germany     & 0.909 & 0.929 & 1.014 & 1.308 & & 6.45 & 5.02 & 4.16 & 1.21 \\
Greece      & 0.929 & 0.946 & 1.111 & 1.541 & & 2.78 & 1.95 & 1.22 & $-$1.38 \\
Spain       & 0.857 & 0.890 & 0.969 & 1.383 & & 2.55 & 2.13 & 1.52 & $-$0.88 \\
France      & 0.926 & 0.948 & 1.041 & 1.140 & & 2.20 & 2.30 & 1.99 & 1.35 \\
Ireland     & 0.816 & 1.052 & 1.574 & 2.167 & & 2.67 & 2.12 & 1.87 & 0.68 \\
Italy       & 0.944 & 0.939 & 0.973 & 1.038 & & 3.09 & 2.52 & 1.77 & 0.65 \\
Netherlands & 0.811 & 0.824 & 0.869 & 1.038 & & 4.99 & 4.22 & 3.09 & 2.12 \\
Portugal    & 1.012 & 1.062 & 1.176 & 1.362 & & 1.77 & 1.17 & 0.43 & $-$0.77 \\
\midrule
Mean        & 0.899 & 0.933 & 1.052 & 1.316 & & & & & \\
Ratio $<1$  & 10 & 9 & 6 & 0 & & & & & \\
CW $>1.645$ & 11 & 10 & 8 & 3 & & & & & \\
\bottomrule
\end{tabular}

\vspace{0.5em}
\begin{minipage}{0.95\textwidth}
\footnotesize \textit{Notes.} Ratio of the root mean squared forecast error of
the index regression to that of the best of the recursive mean, the random walk
and the own-lag regression; values below one favour the index. The random walk
is the best benchmark in every cell except Austria at $h \le 4$, where the
own-lag regression is, and Austria at $h = 8$, where the recursive mean is.
\end{minipage}
\end{table}

Table~\ref{tab:contrast} restates the same forecasts against the own-lag
regression alone. The picture is unrecognisable: the index now wins at every
horizon in almost every region, and the two-year row --- where the honest
comparison shows a loss in eleven of eleven --- shows a win in ten of eleven and
a mean reduction in error of 13\%.

\begin{table}[htbp]
\centering
\caption{The same forecasts, scored two ways}
\label{tab:contrast}
\begin{tabular}{lrrrrcrrrr}
\toprule
& \multicolumn{4}{c}{vs.\ own-lag regression} & & \multicolumn{4}{c}{vs.\ best benchmark} \\
\cmidrule(lr){2-5} \cmidrule(lr){7-10}
& $h{=}1$ & $h{=}2$ & $h{=}4$ & $h{=}8$ & & $h{=}1$ & $h{=}2$ & $h{=}4$ & $h{=}8$ \\
\midrule
Mean ratio      & 0.864 & 0.839 & 0.832 & 0.872 & & 0.899 & 0.933 & 1.052 & 1.316 \\
Ratio $<1$      & 11 & 11 & 11 & 10 & & 10 & 9 & 6 & 0 \\
CW $> 1.645$    & 11 & 10 & 10 & 10 & & 11 & 10 & 8 & 3 \\
\bottomrule
\end{tabular}
\end{table}

\section{Robustness}
\label{sec:robust}

\textbf{Target definition.} Scoring the single quarter $t+h$ rather than the
path over $t+1,\dots,t+h$ makes the index look worse at every horizon beyond
the first: at four quarters it wins in three regions rather than six, and the
mean ratio is 1.106 rather than 1.052; at two years, 1.461 rather than 1.316.
The path target is therefore the conservative choice for the negative result
reported above.

\textbf{Other target variables.} For output growth the recursive mean is the
best benchmark in all eleven regions and the index beats it in none, at any
horizon, with mean ratios between 1.20 and 1.65. For inflation the index wins
in six or seven regions with mean ratios close to one and Clark--West
statistics clearing in at most five. Neither target shows the short-horizon
result that unemployment does.

\textbf{Number of factors.} Extracting a second factor while continuing to use
only the first as a predictor changes the counts in Table~\ref{tab:main} not at
all and the mean ratios by less than a hundredth.

\textbf{A longer sample.} Austria is the one region for which national sources
allow the panel to be extended, to 1993, giving 100 to 107 origins. The
short-horizon result survives, and so does the fact that Austria is atypical:
there the own-lag regression is the best of the three benchmarks at every
horizon, and the index beats it narrowly at two years as well.

\section{Limitations}

The benchmark set is three simple rules. A properly specified univariate model
of each country's unemployment rate would be harder to beat than any of them,
and the short-horizon result should be read as a comparison with a random walk
rather than with the state of the art. The dataset contains no real-time
vintages, so the macro variables are revised figures throughout; this affects
all four forecasts equally but inflates the level of accuracy reported. Finally,
the eleven regions share a monetary policy and much of a business cycle, so the
counts should be read as describing a correlated panel rather than eleven
independent tests.

\section{Conclusion}

A financial conditions index for a euro-area country carries information about
unemployment one and two quarters ahead that survives comparison with a random
walk, in ten and nine of eleven regions. At two years it survives nothing: the
index loses to a random walk in every region, while appearing to win in ten of
eleven when scored against an autoregressive benchmark.

The practical implication is a reporting standard rather than a new method.
Evaluations of financial conditions indices at horizons of a year or more
should include a benchmark that imposes the persistence of the target instead
of estimating it, and should say which benchmark the reported ratios are
against. On this evidence the difference between those two choices is larger
than the difference between having an index and not having one.

\begin{thebibliography}{99}

\bibitem[Barigozzi and Lissona(2024)]{barigozzi2024}
Barigozzi, M., Lissona, C., 2024. EA-MD-QD: Large euro area and euro member
countries datasets for macroeconomic research.
\url{https://github.com/BarigozziMatteo/EA-MD-QD}

\bibitem[Clark and West(2007)]{clark2007}
Clark, T.E., West, K.D., 2007. Approximately normal tests for equal predictive
accuracy in nested models. Journal of Econometrics 138, 291--311.

\bibitem[Hatzius et al.(2010)]{hatzius2010}
Hatzius, J., Hooper, P., Mishkin, F.S., Schoenholtz, K.L., Watson, M.W., 2010.
Financial conditions indexes: a fresh look after the financial crisis. NBER
Working Paper 16150.

\bibitem[Koop and Korobilis(2014)]{koop2014}
Koop, G., Korobilis, D., 2014. A new index of financial conditions. European
Economic Review 71, 101--116.

\end{thebibliography}

\end{document}
